% =======> Project report for CSE 151B Spring 2026 Math Reasoning Competition
% Authors: Connor Wu
% =================================================

\documentclass{article}

% to avoid loading the natbib package, add option nonatbib:
% \usepackage[nonatbib]{neurips_2024}
% \PassOptionsToPackage{square,comma,sort,numbers}{natbib}
\PassOptionsToPackage{numbers, compress}{natbib}
% \usepackage[numbers]{natbib}

% to compile a preprint version, e.g., for submission to arXiv, add add the
% [preprint] option:
\usepackage[preprint]{neurips_2024}
% \usepackage[nonatbib, preprint]{neurips_2024}

\usepackage[utf8]{inputenc} % allow utf-8 input
% \usepackage[T1]{fontenc}    % use 8-bit T1 fonts
% For maths/theorems and such
\usepackage{amsmath}
\usepackage{amssymb}
\usepackage{mathtools}
\usepackage{amsthm}
\usepackage{float}
% Use DeclareMathOperator and/or \newcommand to define abbreviations, e.g.:
\DeclareMathOperator{\R}{\mathbb{R}} % real numbers
% Recommended, but optional, packages for figures and better typesetting:
\usepackage{microtype}
\usepackage{graphicx}
% \usepackage{subfigure}  % removed: conflicts with subcaption on modern texlive
\usepackage{subcaption}
\usepackage{booktabs} % for professional tables (https://www.ctan.org/pkg/booktabs)
\usepackage{algorithm}  % For algorithms
\usepackage{algorithmic}
\usepackage{hyperref} % hyperref makes hyperlinks in the resulting PDF.
\usepackage{wrapfig}  % For wrapping figures with text
\usepackage[capitalize,noabbrev]{cleveref}   % For \cref
\RequirePackage{doi}
% Color references and URLs
\usepackage{xcolor, colortbl}
\definecolor{darkblue}{rgb}{0.19, 0.19, 0.62}
\hypersetup{
  colorlinks = true,
  citecolor  = darkblue,
  filecolor  = darkblue,
  urlcolor   = darkblue,
}
% Todonotes is useful during development; simply uncomment the next line
%    and comment out the line below the next line to turn off comments
%\usepackage[disable,textsize=tiny]{todonotes}
\usepackage[textsize=tiny]{todonotes}

% \input{_macros}



\title{
CSE 151B/251B Milestone Report:\\
Inference-Time Engineering for Math Reasoning with a Quantized Thinking LLM
% \\
% {\small \url{https://github.com/<TODO: project repo URL>}}
}

\author{%
  Connor Wu \\
  Department of Computer Science and Engineering\\
  University of California, San Diego\\
  \texttt{cow003@ucsd.edu} \\
}
\begin{document}
\maketitle


\begin{abstract}
The CSE 151B Spring 2026 competition asks us to maximise accuracy on $1126$
mixed multiple-choice (MCQ) and free-form math problems using a fixed,
INT8-quantised \texttt{Qwen3-4B-Thinking-2507} model. With no training in our
current setup, all gains come from inference-time engineering. We extend the
starter notebook with four targeted changes -- LaTeX repair scoped to math
delimiters, prompts tuned for a thinking model, a lazy-fallback MCQ extractor
with logprobs calibration, and hardened brace-aware boxed-answer parsing --
each with negligible runtime cost. \textbf{Final private Kaggle score:
TODO~--~to be added after the full submission is graded.}
\end{abstract}


% % % === Introduction
\section{Introduction}
\label{sec:intro}

\paragraph{Problem Definition.}
The public split contains $n{=}1126$ items: $375$ MCQs with up to ten labelled
options $A,B,\dots,J$, and $751$ free-form problems whose gold answers are
symbolic or numeric expressions. The model \texttt{Qwen3-4B-Thinking-2507} is
served by vLLM under INT8 BitsAndBytes quantisation; no fine-tuning is part of
the current pipeline. A symbolic judger checks free-form answers for
mathematical equivalence rather than string match.

\paragraph{Problem Significance.}
Quantised inference of mid-sized open-weight reasoning models is the
deployment setting most practitioners actually face: training budgets are out
of reach, but a single consumer GPU can hold a 4--7B model in 8-bit.
Understanding which inference-time choices help -- particularly for
``thinking'' models that emit a long
$\langle\texttt{think}\rangle\dots\langle/\texttt{think}\rangle$ chain before
the answer -- transfers directly to real low-compute applications.

\paragraph{Technical Challenge.}
Three challenges shape the project. (1)~\textit{Data quality:} the dataset
stores LaTeX with backslashes stripped (\texttt{frac}, \texttt{infty},
\texttt{int}, etc.), and a naive repair corrupts common English words
(``sum'', ``times'', ``log'') that share spellings with LaTeX commands.
(2)~\textit{Thinking-model output handling:} the model emits a long internal
chain before the final boxed answer; the token budget and stop conditions
must accommodate it. (3)~\textit{Robust scoring:} MCQ scoring is letter
match, but free-form scoring requires the response to contain exactly one
well-formed \texttt{\textbackslash boxed\{\}} that survives nested-brace
parsing.

\paragraph{Contributions.}
The starter code is a working end-to-end pipeline (vLLM, $32$k-token budget,
simple system prompts, regex letter extraction, symbolic judger). Its
limitations: no LaTeX repair, MCQ extraction never abstains (it falls back to
the last capital letter anywhere in the response), and \texttt{\textbackslash
boxed\{\}} contents with nested braces are truncated. Our changes:
\begin{itemize}
  \item \textbf{Scoped LaTeX repair} that restores backslashes only inside
    \texttt{\$\ldots\$} math blocks, leaving English untouched.
  \item \textbf{Thinking-tuned prompts} -- the MCQ prompt explicitly directs
    the model to emit \texttt{\textbackslash boxed\{X\}} after
    \texttt{</think>}, and the free-form prompt includes two in-context
    examples enforcing a single boxed answer with comma-separated
    sub-answers.
  \item \textbf{Lazy-fallback MCQ extractor} -- the thinking response is the
    primary source; a cheap no-think logprobs pass over letter tokens is
    invoked only when the primary extractor abstains.
  \item \textbf{Hardened answer parsing} -- \texttt{extract\_letter} reads
    preferentially from the post-\texttt{</think>} region, and
    \texttt{extract\_boxed} uses a brace-counting parser so that
    \texttt{\textbackslash boxed\{\textbackslash frac\{1\}\{2\}\}} extracts
    in full.
\end{itemize}


% % % === Related Work (kept brief; competition allows but does not require this)
\section{Related Work}
\label{sec:relatedwork}

\paragraph{Reasoning LLMs.} Recent open-weight ``thinking'' models such as the
Qwen3 series~\citep{qwen3} interleave a hidden chain of thought, delimited by
\texttt{<think>}/\texttt{</think>}, with a final user-facing answer. Their
accuracy on math benchmarks improves with adequate decoding budget and
suitable sampling temperature; recommended defaults are documented by the
model authors.

\paragraph{Inference serving.} vLLM~\citep{vllm} provides batched, paged
attention serving of HuggingFace models, including BitsAndBytes 8-bit weight
quantisation. Decoding throughput depends on \texttt{max\_num\_seqs},
\texttt{enable\_prefix\_caching}, and whether CUDA graphs are enabled
(\texttt{enforce\_eager=False}).

\paragraph{Symbolic answer judging.} Math benchmarks such as
MATH~\citep{hendrycks2021math} have driven the design of symbolic equivalence
checkers built on SymPy-like CAS backends; the course-provided
\texttt{Judger} class follows this lineage and accepts free-form responses,
extracting and normalising the boxed expression internally.


% % % === Methods
\section{Methods}
\label{sec:methods}

\paragraph{Problem Setting.}
Let $\mathcal{D}=\{(q_i,O_i,y_i)\}_{i=1}^{n}$ be the public set with $O_i$
possibly empty. The model $f_\theta$ is held fixed; we choose a prompt template
$\pi$, sampling parameters $\sigma$, and extraction function $g$ to maximise
\[
  \mathrm{Acc}(\pi,\sigma,g) \;=\; \tfrac{1}{n}\sum_{i=1}^{n}\,
    \mathbf{1}\!\left[\mathrm{Judge}\bigl(g(f_\theta(\pi(q_i,O_i);\sigma))\bigr) = y_i\right],
\]
with $\mathrm{Judge}$ being letter equality for MCQ and symbolic equivalence
for free-form. There is no trainable loss in this milestone; all optimisation
is over the discrete choices $(\pi,\sigma,g)$.

\paragraph{Changes from the starter.}
\Cref{tab:changes} summarises the differences between the starter and our
current optimised notebook. Each change is small and local, and the
lazy-fallback design means the optimised version costs no more at decode time
than the starter.

\begin{table}[h]
  \caption{Differences between the starter and the current optimised
    notebook. All changes are inference-time only; the model is unchanged.}
  \label{tab:changes}
  \centering
  \footnotesize
  \begin{tabular}{lp{4.6cm}p{4.6cm}}
    \toprule
    Component & Starter & Optimised \\
    \midrule
    Input preprocessing & None. & LaTeX repair scoped to \texttt{\$\ldots\$}
      blocks: restores backslashes on \texttt{frac, infty, int, sum, \ldots}
      inside math while leaving English untouched. \\[2pt]
    MCQ prompt & Generic ``output the letter in \texttt{\textbackslash
      boxed\{\}}.'' & Tuned for a thinking model: ``After \texttt{</think>},
      output exactly one line: \texttt{\textbackslash boxed\{X\}}.'' \\[2pt]
    Free-form prompt & Generic step-by-step instruction. & Adds two
      in-context examples enforcing a single \texttt{\textbackslash
      boxed\{\}} with comma-separated sub-answers. \\[2pt]
    MCQ extraction & Regex that falls back to the last capital letter
      anywhere in the response (never abstains). & Multi-pattern extractor
      that prefers text after \texttt{</think>} and abstains otherwise; on
      abstention, a cheap no-think logprobs pass reads $P(\text{letter})$
      over $A$--$J$ tokens. \\[2pt]
    Free-form parsing & Raw response sent to the judger. & Same -- plus a
      brace-counting \texttt{extract\_boxed} used for diagnostics. \\[2pt]
    Sampling / budget & $T{=}0.6$, $p{=}0.95$, $k{=}20$, \texttt{max\_tokens}
      $=32768$, \texttt{enable\_prefix\_caching=False}. & Same sampling;
      \texttt{max\_tokens}$=16384$ (ample in practice, lighter on the KV
      cache); \texttt{enable\_prefix\_caching=True}. \\
    \bottomrule
  \end{tabular}
\end{table}


% % % === Experiments
\section{Experiments}
\label{sec:experiments}

\subsection{Baselines}
\label{sec:baselines}

We compare three configurations on the public development data:
\begin{itemize}
    \item \textbf{Starter} -- the unmodified course-provided notebook.
    \item \textbf{Optimised} -- this milestone, with the changes summarised
      in \cref{tab:changes}.
\end{itemize}

\subsection{Evaluation}
\label{sec:evaluation}

MCQ predictions are scored by case-insensitive letter equality with the gold
label. Free-form predictions are scored by the course-provided symbolic
\texttt{Judger}, which extracts and normalises the boxed expression and
tests mathematical equivalence against each gold answer. We report MCQ,
free-form, and overall accuracy.

\subsection{Implementation details}
\label{sec:implementation_details}

The model is \texttt{Qwen/Qwen3-4B-Thinking-2507} served by vLLM $0.20.2$
with \texttt{quantization="bitsandbytes"}, \texttt{max\_model\_len}=16384,
\texttt{gpu\_memory\_utilization}=$0.50$. Sampling follows the Qwen3-Thinking
recommendation ($T{=}0.6$, $p{=}0.95$, $\mathrm{top\_k}{=}20$,
\texttt{repetition\_penalty}$=1.0$), with \texttt{max\_tokens}$=16{,}384$ on
the main pass and a near-greedy fallback ($T{=}0.01$, $p{=}1.0$,
\texttt{logprobs}=20) for MCQ logprobs. GPU footprint at load is
$\approx$2.7~GiB of weights plus an $8.3$~GiB KV cache, fitting in
$\sim$16~GiB of VRAM. \todo{Add public GitHub URL for final report.}

\subsection{Results}
\label{sec:results}

\Cref{tab:dev-results} reports the starter's preliminary $N{=}5$ smoke-test
accuracy (its default), recorded directly from the notebook output. A full
$N{=}1126$ run of the optimised notebook is in progress; final numbers will
replace the placeholder row in the camera-ready version of this report.

\begin{table}[h]
  \caption{Preliminary accuracy on the public set. The starter cell runs
    only the first $5$ items by default. Optimised numbers are pending the
    full-dataset run.}
  \label{tab:dev-results}
  \centering
  \begin{tabular}{lccc}
    \toprule
    Setting & MCQ & Free-form & Overall \\
    \midrule
    Starter ($N{=}5$)   & $1/2$ (50.0\%) & $2/3$ (66.7\%) & $3/5$ (60.0\%) \\
    Optimised ($N{=}1126$) & \multicolumn{3}{c}{TODO -- full-dataset run pending} \\
    \bottomrule
  \end{tabular}
\end{table}

The optimised configuration is no more expensive than the starter at decode
time: \texttt{max\_tokens} is lower, the fallback calibration pass is
invoked only for MCQs where the primary extractor abstains, and
\texttt{enable\_prefix\_caching=True} allows the system prompt to be reused
across the batch.


% % % === Discussion
\section{Discussion}
\label{sec:discussion}

Thinking models are sensitive to anything that prevents the final answer from
surfacing: token budgets that truncate the chain, scoring paths that prefer
a no-think pass, and \texttt{stop=[\ldots]} lists that can fire during
reasoning all act as silent answer-suppressors. The optimised notebook
addresses these failure modes upfront.

\paragraph{Achievements.}
A small, runtime-flat delta to the starter that adds scoped LaTeX repair,
thinking-tuned prompts, a lazy-fallback MCQ extractor with logprobs
calibration, and brace-aware boxed parsing.

\paragraph{Bottlenecks and Mitigation.}
The dominant bottleneck is decode latency: a full $1126$-question run takes
a non-trivial fraction of an hour, which makes hyperparameter sweeps
expensive. We plan to maintain a frozen $\sim$100-item stratified subsample
(75 free-form, 25 MCQ) as a stable evaluation surface for fast iteration.

\paragraph{Plans Forward.}
Three directions are queued: (i)~self-consistency over $n{=}5$ free-form
samples scored by the symbolic judger; (ii)~a cheap no-think first pass used
to \emph{filter} which items deserve the expensive thinking pass on rerun;
(iii)~a post-processor that normalises boxed expressions to make the judger
robust on borderline cases.

\paragraph{Team Responsibilities.}
Solo submission. Connor Wu is responsible for the full pipeline, evaluation
harness, and report.

\paragraph{Timelines.}
\begin{itemize}
    \item Weeks 1--3: environment setup, starter notebook end-to-end,
      dataset and judger inspection. \emph{(Complete.)}
    \item Weeks 4--5: optimised notebook (this milestone). \emph{(Complete.)}
    \item Week 6: full $N{=}1126$ evaluation and first leaderboard submission.
    \item Week 7: self-consistency / majority voting on free-form.
    \item Week 8: cheap-first-pass filtering of items for the thinking pass.
    \item Week 9: post-processing pipeline (numeric normalisation).
    \item Week 10: final ablations and submission.
      \todo{Adjust once final-quarter dates are confirmed.}
\end{itemize}


\section*{LLM Usage}

Anthropic's Claude (model \texttt{claude-opus-4-7}, via the Claude desktop
``Cowork'' interface) was used as an interactive coding assistant with
read/write access to this repository. Claude reviewed the diff between the
starter and the optimised notebook, proposed and applied the changes
summarised in \cref{tab:changes}, and drafted this report from a short
briefing on scope, authorship, and which numbers to cite. The author ran all
experiments, verified the patched notebook, and reviewed and revised every
claim in the report. The author takes full responsibility for the contents.


% % % === ACKNOWLEDGMENTS AND REFERENCES START
% \section*{Acknowledgements}
% This work was supported $\dots$.
\bibliography{references.bib}
\bibliographystyle{plainnat}  % Can also use abbrvnat, unsrt, etc.

% % % === APPENDIX (optional)
% \newpage
% \appendix
% \section*{Appendix}
% \section{Appendix topic 1}
%
\end{document}
